
\section{Conversion changes a claim and the people who control it}
\label{sim:conversion-book-rights}
Preserving rights across a reporting boundary is necessary until an event
changes them. A conversion provides a contrasting case: cash may stay fixed
while debt disappears, new shares are issued and another holder gains
control. A simulator that checks cash and balanced journals alone can miss
the decision consequence.

Consider a separately supplied fiscal opening on 1 January 2028. Project
assets are 100 and cash is zero. Senior debt has face and carrying amount
40. A convertible claim has face 80 and a supplied carrying amount $C$.
There are 40 old shares and no shares held by the convertible creditor.
Opening capital is $60-C$, so the signed balances reconcile. The examples
use $C=80$ and $C=17000/441$ as two distinct accounting premises. The latter
is a diagnostic amount; selecting it does not establish a measurement rule.

The supplied contract says that the conversion signal extinguishes the
claim and issues 120 shares to its holder on that date. Under the selected
carrying-value transfer, the journal debits the convertible liability by
$C$ and credits issued equity by $C$. The face becomes zero, new shares
become 120 and the observed signal and conversion are recorded. There is
no cash flow or profit from this assumed transfer. The financial statements
then show assets 100, liabilities 40 and equity 60 in either carrying case.
Under the other signal, the debt remains outstanding and no shares issue.

Cash-flow and voting fractions follow from the resulting shares under the
assumed equal rights per share. After conversion the new holder owns
$120/(40+120)=3/4$ and the old holder owns $1/4$. The new holder therefore
has strict majority control. Before conversion the old holder controls all
shares. These rights follow from the completed transition, rather than
from a controller label retained from the opening. Different voting classes
would require a different supplied rights rule.

The complete comparisons retain every account, the claim and share record,
unrelated memory, event history, journal and report on two account charts.
They also compare the resulting rights with a separately authored payoff
table. Omitting claim extinction or share issue fails even when a proposed
journal balances. Repeating conversion with a fresh event identifier fails
because the event has already occurred. Wrong dates and signals fail before
effects are accepted. Reporting the equity increase without explaining its
noncash movement also fails, while preserving the already accepted conversion.

This example starts from an explicit opening and assumes the accounting
transfer. It does not derive that opening from issuance, select a reporting
standard, or establish that a real contract has these terms. Its contribution
is narrower and useful: the accounting balances and economic rights change
together, and the rights can supply the next decision without counting both
converted debt and equity as surviving claims.
